\documentclass[11pt, a4paper]{article}
\usepackage[margin=2.5cm]{geometry}
\usepackage{booktabs}
\usepackage{array}
\usepackage{longtable}
\usepackage{xcolor}
\usepackage{hyperref}
\usepackage{microtype}
\usepackage{parskip}
\usepackage{enumitem}
\usepackage[T1]{fontenc}
\usepackage[utf8]{inputenc}

\definecolor{highprio}{HTML}{CC0000}
\definecolor{medprio}{HTML}{CC8800}
\definecolor{lowprio}{HTML}{007700}

\hypersetup{
  colorlinks = true,
  linkcolor  = blue!60!black,
  urlcolor   = blue!60!black,
}

\title{\textbf{ERBOT Results Summary}\\[4pt]
       \large Synthetic 10K Dataset}
\author{ERBOT v0.2.0}
\date{Generated: 2026-03-23 \quad|\quad Run time: $\approx$11 minutes (SVM excluded)}

\begin{document}
\maketitle
\tableofcontents
\vspace{1em}
\hrule
\vspace{1em}

% ─────────────────────────────────────────────────────────────────────────────
\section{Dataset Overview}

\begin{table}[h!]
\centering
\begin{tabular}{ll}
\toprule
\textbf{Property} & \textbf{Value} \\
\midrule
Dataset      & Synthetic 10K (d10k) \\
Records      & 10,000 \\
Fields       & 2 (\texttt{id}, \texttt{text}) \\
Task         & Deduplication (single-source) \\
Ground truth & 8,705 matched pairs (from \texttt{10Kduplicates.csv}) \\
\bottomrule
\end{tabular}
\end{table}

The synthetic 10K dataset consists of 10,000 noisy text records generated to simulate
real-world deduplication challenges. Each record has a single free-text field
(\texttt{text}) containing a scrambled mix of tokens --- names, dates, numbers,
addresses --- with no explicit field structure. Duplicates are introduced by corrupting
original records with character substitutions, token deletions, and reorderings,
mimicking common data entry and OCR errors.

% ─────────────────────────────────────────────────────────────────────────────
\section{Pipeline Configuration}

\begin{table}[h!]
\centering
\begin{tabular}{lll}
\toprule
\textbf{Stage} & \textbf{Setting} & \textbf{Value} \\
\midrule
Blocking   & Method          & Prefix blocking on \texttt{text}, length 3 \\
Blocking   & Candidate pairs & 885,301 \\
Blocking   & Reduction ratio & 0.982 (98.2\% of all pairs eliminated) \\
Similarity & Fields computed & 1 (\texttt{text} only) \\
Weights    & Method          & Equal (weight = 1.0 for \texttt{text}) \\
Clustering & Methods run     & 8 (SVM excluded --- see note) \\
Merge      & Strategy        & Best \\
\bottomrule
\end{tabular}
\end{table}

\textbf{SVM excluded:} SVM was initially run but terminated after $\approx$20 hours
without completing. With 885,301 candidate pairs, SVM training time was estimated at
1--4 days due to its super-linear scaling with the number of pairs. All other 8 methods
completed in $\approx$11 minutes. SVM remains viable for smaller datasets (CORA at
137K pairs completed in $\approx$63 minutes total) but is not practical at this scale
without subsampling or approximation. A scalable alternative --- logistic regression or
gradient boosting --- is recommended as a replacement.

\textbf{Blocking note:} 885,301 pairs from 10,000 records (49,995,000 total possible).
Prefix-3 achieves 98.2\% reduction --- the highest across all three datasets tested.
However, for a single noisy text field, many true matches will have corrupted or
reordered opening tokens and are silently discarded at this stage.

% ─────────────────────────────────────────────────────────────────────────────
\section{Performance Results}

8 clustering methods evaluated (SVM excluded):

\begin{table}[h!]
\centering
\small
\setlength{\tabcolsep}{4pt}
\begin{tabular}{lrrrrrrrrrr}
\toprule
\textbf{Method} & \textbf{ARI} & \textbf{NMI} & \textbf{VI} &
\textbf{B$^3$-P} & \textbf{B$^3$-R} & \textbf{B$^3$-F} &
\textbf{Hom.} & \textbf{Com.} & \textbf{V} \\
\midrule
leiden           & 0.000 & \textbf{0.914} & \textbf{1.97} & \textbf{1.000} & 0.271 & \textbf{0.426} & \textbf{1.000} & 0.841 & \textbf{0.914} \\
\textbf{threshold\_cc} & \textbf{0.023} & 0.778 & 4.21 & 0.334 & 0.519 & 0.406 & 0.707 & 0.865 & 0.778 \\
\textbf{louvain} & \textbf{0.023} & 0.778 & 4.21 & 0.334 & 0.519 & 0.406 & 0.707 & 0.865 & 0.778 \\
\textbf{label\_prop} & \textbf{0.023} & 0.778 & 4.21 & 0.334 & 0.519 & 0.406 & 0.707 & 0.865 & 0.778 \\
hclust\_avg      & 0.002 & 0.371 & 8.82 & 0.019 & 0.573 & 0.036 & 0.249 & 0.727 & 0.371 \\
pam              & 0.002 & 0.371 & 8.82 & 0.019 & 0.573 & 0.036 & 0.249 & 0.727 & 0.371 \\
hclust\_ward     & 0.001 & 0.347 & 8.80 & 0.019 & 0.690 & 0.037 & 0.224 & 0.772 & 0.347 \\
gc               & 0.000 & 0.000 & 10.4 & 0.001 & 1.000 & 0.002 & 0.000 & 1.000 & 0.000 \\
\midrule
\textit{final}   & 0.023 & 0.778 & 4.21 & 0.334 & 0.519 & 0.406 & 0.707 & 0.865 & 0.778 \\
\bottomrule
\end{tabular}
\end{table}

\textit{Metrics: ARI = Adjusted Rand Index; NMI = Normalized Mutual Information;
VI = Variation of Information (lower is better); B$^3$ = B-cubed; V = V-measure;
Hom.\ = Homogeneity; Com.\ = Completeness. Full PairF values in \texttt{er\_performance.csv}.}

% ─────────────────────────────────────────────────────────────────────────────
\section{Final Cluster Assignment}

The \texttt{final} strategy correctly selected \textbf{threshold\_cc} (highest ARI = 0.023).

\begin{table}[h!]
\centering
\begin{tabular}{lr}
\toprule
\textbf{Statistic} & \textbf{Value} \\
\midrule
Total clusters    & 1,619 \\
Records           & 10,000 \\
Mean cluster size & 6.2 \\
\bottomrule
\end{tabular}
\end{table}

% ─────────────────────────────────────────────────────────────────────────────
\section{Method-by-Method Interpretation}

\subsection*{Threshold-CC / Louvain / Label Propagation (Highest ARI --- 0.023)}
The three graph-based methods again converge to identical solutions, confirming that
similarity graph structure dominates algorithm choice. B$^3$-F of 0.406 with precision
0.334 and recall 0.519 is a moderate result --- about half of true matches are found at
the cost of some false merges. These are the best practical methods at this scale.

\subsection*{Leiden (ARI = 0.000, NMI = 0.914, B$^3$-F = 0.426)}
Leiden produces the highest NMI and B$^3$-F but zero ARI. This apparent contradiction
arises because ARI heavily penalises over-splitting: Leiden creates many small, highly
pure sub-clusters (B$^3$-precision = 1.000) but each true entity is fragmented across
multiple predicted clusters (recall = 0.271). The clusters are perfectly pure but
incomplete.

For applications where \textbf{false merges are very costly} (e.g.\ legal identity
resolution where merging two different people is a serious error), Leiden at its default
resolution is actually the safest choice. For applications where \textbf{recall matters},
threshold-CC is better.

\subsection*{hclust\_avg / PAM (ARI $\approx$ 0.002)}
Near-zero ARI despite moderate NMI. Both achieve near-zero precision (0.019), meaning
almost every pair they merge is a false positive. At 10K records the similarity matrix
is too sparse and noisy for centroidal/hierarchical methods to find meaningful structure.

\subsection*{hclust\_ward (ARI = 0.001)}
Similar to hclust\_avg. Ward linkage performs marginally worse on precision but slightly
better on recall. Hierarchical methods are clearly not suited to large sparse similarity
matrices.

\subsection*{GC --- Graph Coloring (Degenerate --- ARI = 0.000)}
Perfect recall, near-zero precision, collapsed to one giant cluster. The gc dimension
mismatch error persists across all three datasets --- this is a recurring bug in the
GCMER interface layer that needs fixing.

\subsection*{SVM (Not Run)}
SVM was launched but killed after $\approx$20 hours still processing 885,301 pairs.
Based on observed runtimes on smaller datasets (CORA 137K pairs: included in 63 min;
affiliation 164K pairs: dominated a 5.6-hour run), estimated completion time for 885K
pairs was 1--4 days. SVM's kernel matrix computation scales super-linearly with the
number of training examples, making it impractical at this dataset size. Logistic
regression or gradient boosting are recommended as drop-in replacements with
near-linear scaling.

% ─────────────────────────────────────────────────────────────────────────────
\section{Key Findings}

\begin{enumerate}[leftmargin=*, label=\arabic*.]
  \item \textbf{Overall performance is very low.} Best ARI is 0.023 --- lower than
        affiliation (0.076) and far below CORA (0.746). At 10,000 records with a single
        noisy unstructured text field, the task is fundamentally harder.

  \item \textbf{Blocking discards too many true matches.} 98.2\% reduction is
        aggressive. With only 8,705 gold pairs in a 10,000-record dataset, many duplicates
        involve records with corrupted or reordered first tokens, which prefix-3 blocking
        eliminates silently. Blocking is the primary performance bottleneck.

  \item \textbf{Leiden is the precision champion.} Perfect B$^3$-precision (1.000) and
        NMI (0.914) make Leiden the best choice when false merges must be avoided.
        Its ARI = 0 reflects over-splitting, not noise.

  \item \textbf{SVM is not viable at this scale.} For datasets above $\approx$200K
        candidate pairs, SVM should be replaced with a faster supervised classifier.

  \item \textbf{Graph methods plateau quickly.} Threshold-CC, Louvain, and Label
        Propagation give identical results here as in the previous two datasets. The
        similarity graph structure consistently dominates community detection algorithm
        choice.

  \item \textbf{The \texttt{final} merge correctly selected threshold\_cc.} Merge bug
        fix continues to work correctly across all three datasets.
\end{enumerate}

% ─────────────────────────────────────────────────────────────────────────────
\section{Recommendations for Next Steps}

\begin{table}[h!]
\centering
\begin{tabular}{cp{0.78\linewidth}}
\toprule
\textbf{Priority} & \textbf{Action} \\
\midrule
{\color{highprio}\textbf{High}}  & Replace prefix-3 blocking with SN blocking (window = 40--80) --- the 98.2\% reduction is too aggressive for a noisy single-field dataset. \\[4pt]
{\color{highprio}\textbf{High}}  & Replace SVM with logistic regression or gradient boosting for supervised classification --- near-linear scalability, similar accuracy. \\[4pt]
{\color{medprio}\textbf{Medium}} & Tune Leiden resolution downward (try 0.1--0.3) to improve recall while preserving its high precision. \\[4pt]
{\color{medprio}\textbf{Medium}} & Fix GC dimension mismatch error --- silently falls back to degenerate solution across all datasets. \\[4pt]
{\color{medprio}\textbf{Medium}} & Measure pair completeness at Stage 3 to quantify how many gold pairs are lost at blocking. \\[4pt]
{\color{lowprio}\textbf{Low}}    & Explore token-set / Jaccard similarity instead of default string similarity for unstructured text. \\[4pt]
{\color{lowprio}\textbf{Low}}    & Try field parsing: split \texttt{text} into sub-fields (name, date, address tokens) to enable multi-field similarity. \\
\bottomrule
\end{tabular}
\end{table}

% ─────────────────────────────────────────────────────────────────────────────
\section{Cross-Dataset Comparison}

\begin{table}[h!]
\centering
\begin{tabular}{llll}
\toprule
\textbf{Metric} & \textbf{CORA} & \textbf{Affiliation} & \textbf{Syn10k} \\
\midrule
Records              & 1,879   & 2,260   & 10,000 \\
Fields (usable)      & 15      & 1       & 1 \\
Gold pairs           & 64,578  & 32,816  & 8,705 \\
Candidate pairs      & 137,707 & 164,719 & 885,301 \\
Blocking RR          & 92.2\%  & 93.5\%  & 98.2\% \\
Best ARI             & 0.746   & 0.076   & 0.023 \\
Best B$^3$-F         & 0.738   & 0.460   & 0.426 \\
Best method (ARI)    & hclust\_avg & threshold\_cc & threshold\_cc \\
SVM viable?          & Yes ($\approx$63 min) & Marginal ($\approx$5.6h) & No ($\approx$1--4 days) \\
Final correct?       & Yes     & Yes     & Yes \\
\bottomrule
\end{tabular}
\end{table}

The clear trend: \textbf{more records + fewer fields = lower performance}. CORA benefits
from 15 discriminative fields that collectively identify duplicates even under noise.
Affiliation and Syn10k each have one field, forcing all discriminative power through a
single noisy similarity score.

% ─────────────────────────────────────────────────────────────────────────────
\section{Output Files}

\begin{table}[h!]
\centering
\begin{tabular}{ll}
\toprule
\textbf{File} & \textbf{Description} \\
\midrule
\texttt{er\_performance.csv}          & Full metric table (9 methods $\times$ 13 metrics) \\
\texttt{er\_predictions.csv}          & Final cluster assignments (10,000 records $\times$ 2 columns) \\
\texttt{d10k\_data\_profile.pdf}      & Dataset profile: field types, distributions, missingness \\
\texttt{plot\_method\_comparison.pdf} & Grouped bar chart: ARI, B$^3$-F, NMI per method \\
\texttt{plot\_performance\_heatmap.pdf} & Colour heatmap: all metrics $\times$ all methods \\
\texttt{plot\_cluster\_sizes.pdf}     & Cluster size distribution histogram \\
\texttt{plot\_precision\_recall.pdf}  & Precision vs.\ Recall scatter (B-cubed and Pair-F) \\
\texttt{SYN10K\_RESULTS\_SUMMARY.pdf} & This document \\
\bottomrule
\end{tabular}
\end{table}

\end{document}
